\documentclass[10pt,a4paper]{amsart}
\usepackage{amsmath,amssymb,amsthm,graphicx}
% Retained minimal preamble for the Stacks Project chapter flat.tex (More on Flatness).
% Verbatim excerpt of the chapter preamble: Xy-pic setup (native preamble lines 14-18) and the native
% theorem declarations with the mathematical macros (native preamble lines 154-217). Copied mechanically;
% nothing added, removed or redefined. The wrapper class and the amsmath/amssymb/amsthm/graphicx base
% packages are supplied by the assembler, so the class/externaldocument part of the native preamble is omitted.
% For commutative diagrams we use Xy-pic
\usepackage[all]{xy}

% We use 2cell for 2-commutative diagrams.
\xyoption{2cell}
% Theorem environments.
%
\theoremstyle{plain}
\newtheorem{theorem}[subsection]{Theorem}
\newtheorem{proposition}[subsection]{Proposition}
\newtheorem{lemma}[subsection]{Lemma}

\theoremstyle{definition}
\newtheorem{definition}[subsection]{Definition}
\newtheorem{example}[subsection]{Example}
\newtheorem{exercise}[subsection]{Exercise}
\newtheorem{situation}[subsection]{Situation}

\theoremstyle{remark}
\newtheorem{remark}[subsection]{Remark}
\newtheorem{remarks}[subsection]{Remarks}

\numberwithin{equation}{subsection}

% Macros
%
\def\lim{\mathop{\mathrm{lim}}\nolimits}
\def\colim{\mathop{\mathrm{colim}}\nolimits}
\def\Spec{\mathop{\mathrm{Spec}}}
\def\Hom{\mathop{\mathrm{Hom}}\nolimits}
\def\Ext{\mathop{\mathrm{Ext}}\nolimits}
\def\SheafHom{\mathop{\mathcal{H}\!\mathit{om}}\nolimits}
\def\SheafExt{\mathop{\mathcal{E}\!\mathit{xt}}\nolimits}
\def\Sch{\mathit{Sch}}
\def\Mor{\mathop{\mathrm{Mor}}\nolimits}
\def\Ob{\mathop{\mathrm{Ob}}\nolimits}
\def\Sh{\mathop{\mathit{Sh}}\nolimits}
\def\NL{\mathop{N\!L}\nolimits}
\def\CH{\mathop{\mathrm{CH}}\nolimits}
\def\proetale{{pro\text{-}\acute{e}tale}}
\def\etale{{\acute{e}tale}}
\def\QCoh{\mathit{QCoh}}
\def\Ker{\mathop{\mathrm{Ker}}}
\def\Im{\mathop{\mathrm{Im}}}
\def\Coker{\mathop{\mathrm{Coker}}}
\def\Coim{\mathop{\mathrm{Coim}}}

% Boxtimes
%
\DeclareMathSymbol{\boxtimes}{\mathbin}{AMSa}{"02}

%
% Macros for moduli stacks/spaces
%
\def\QCohstack{\mathcal{QC}\!\mathit{oh}}
\def\Cohstack{\mathcal{C}\!\mathit{oh}}
\def\Spacesstack{\mathcal{S}\!\mathit{paces}}
\def\Quotfunctor{\mathrm{Quot}}
\def\Hilbfunctor{\mathrm{Hilb}}
\def\Curvesstack{\mathcal{C}\!\mathit{urves}}
\def\Polarizedstack{\mathcal{P}\!\mathit{olarized}}
\def\Complexesstack{\mathcal{C}\!\mathit{omplexes}}
% \Pic is the operator that assigns to X its picard group, usage \Pic(X)
% \Picardstack_{X/B} denotes the Picard stack of X over B
% \Picardfunctor_{X/B} denotes the Picard functor of X over B
\def\Pic{\mathop{\mathrm{Pic}}\nolimits}
\def\Picardstack{\mathcal{P}\!\mathit{ic}}
\def\Picardfunctor{\mathrm{Pic}}
\def\Deformationcategory{\mathcal{D}\!\mathit{ef}}

% --- appended 2026-10-12 (Chief): the project own definitions of the reference/slogan/history
% display environments, copied mechanically from the shared Stacks preamble lines 8-12
% (\usepackage{verbatim} plus the three \newenvironment lines, which the chapter preamble keeps
% verbatim from the project preamble). Nothing added, removed or redefined: these environments
% carry no mathematical content and the project itself typesets them as comments. Without these
% three lines a counted statement whose own text uses them fails with Environment slogan undefined.
\usepackage{verbatim}
\newenvironment{reference}{\comment}{\endcomment}
\newenvironment{slogan}{\comment}{\endcomment}
\newenvironment{history}{\comment}{\endcomment}
\usepackage{cleveref}
\crefname{originallemma}{original lemma}{original lemmas}
\Crefname{originallemma}{Original lemma}{Original lemmas}
\makeatletter
\newlabel{definition-pure}{{05J4}{source}{}{sourcepointer-1}{}}
\newlabel{definition-pure@cref}{{[originallemma][0][]05J4}{[1][0][]source}{}{}{}}
\makeatother
\crefname{originallemma}{original lemma}{original lemmas}
\Crefname{originallemma}{Original lemma}{Original lemmas}
\makeatletter
\newlabel{divisors-lemma-base-change-relative-assassin}{{05DC}{source}{}{sourcepointer-2}{}}
\newlabel{divisors-lemma-base-change-relative-assassin@cref}{{[originallemma][0][]05DC}{[1][0][]source}{}{}{}}
\makeatother
\crefname{originallemma}{original lemma}{original lemmas}
\Crefname{originallemma}{Original lemma}{Original lemmas}
\makeatletter
\newlabel{divisors-remark-base-change-relative-assassin}{{05KL}{source}{}{sourcepointer-3}{}}
\newlabel{divisors-remark-base-change-relative-assassin@cref}{{[originallemma][0][]05KL}{[1][0][]source}{}{}{}}
\makeatother
\crefname{originallemma}{original lemma}{original lemmas}
\Crefname{originallemma}{Original lemma}{Original lemmas}
\makeatletter
\newlabel{lemma-associated-point-specializes}{{05L3}{source}{}{sourcepointer-4}{}}
\newlabel{lemma-associated-point-specializes@cref}{{[originallemma][0][]05L3}{[1][0][]source}{}{}{}}
\makeatother
\crefname{originallemma}{original lemma}{original lemmas}
\Crefname{originallemma}{Original lemma}{Original lemmas}
\makeatletter
\newlabel{lemma-etale-weak-assassin-up-down}{{05FP}{source}{}{sourcepointer-5}{}}
\newlabel{lemma-etale-weak-assassin-up-down@cref}{{[originallemma][0][]05FP}{[1][0][]source}{}{}{}}
\makeatother
\crefname{originallemma}{original lemma}{original lemmas}
\Crefname{originallemma}{Original lemma}{Original lemmas}
\makeatletter
\newlabel{lemma-flat-descend-pure}{{05J9}{source}{}{sourcepointer-6}{}}
\newlabel{lemma-flat-descend-pure@cref}{{[originallemma][0][]05J9}{[1][0][]source}{}{}{}}
\makeatother
\crefname{originallemma}{original lemma}{original lemmas}
\Crefname{originallemma}{Original lemma}{Original lemmas}
\makeatletter
\newlabel{lemma-quasi-finite-base-change}{{05J7}{source}{}{sourcepointer-7}{}}
\newlabel{lemma-quasi-finite-base-change@cref}{{[originallemma][0][]05J7}{[1][0][]source}{}{}{}}
\makeatother
\title{A criterion for purity along a fibre in terms of the relative associated primes inside the image of an etale neighbourhood}
\author{The Stacks Project Authors}
\date{stacks-project:flat.tex:local-snapshot-20261011; GNU FDL 1.2 or later}
\begin{document}
\maketitle
\noindent\textit{Editorial note.} All delivered mathematical text is the author's own, copied from the pinned Stacks Project chapter file flat.tex; no statement or proof step was rewritten, reordered or removed.\par\medskip
\noindent\textit{Editorial note.} Preamble provenance: the chapter's own retained minimal preamble (sources/stacks-flat/retained-preamble.tex, lines 1-85) is a verbatim excerpt of the native Xy-pic setup and of the native theorem declarations and mathematical macros, followed by the project's own reference/slogan/history display-environment definitions copied mechanically from the shared preamble; the wrapper class and the amsmath/amssymb/amsthm/graphicx base packages are supplied by the assembler, and extra\_packages is left empty because everything needed is already in the retained preamble. No macro was invented or redefined.\par\medskip
\noindent\textit{Editorial note.} Same-chapter pointers are located at their \textbackslash{}label lines inside the pinned chapter file; cross-chapter pointers are located in the project's own tag list staged in the snapshot (sources/stacks-chow/tags.txt), at the verbatim lines carrying their canonical tags.\par\medskip
\noindent\textit{Allowed standard prerequisites (editorial).}\begin{itemize}
\item Purity of the base change along an etale neighbourhood and flat descent of purity. Same chapter flat.tex, labels lemma-quasi-finite-base-change (proof line 5226) and lemma-flat-descend-pure (proof line 5228). Shows that purity of F along the fibre is equivalent to purity of the pullback along the elementary etale neighbourhood, which is what allows the base to be replaced by S' and then by its localisation.
\item Associated points of the relative situation and specialization of associated points to the fibre. Same chapter flat.tex, label lemma-associated-point-specializes (proof line 5239); Divisors, label divisors-lemma-base-change-relative-assassin (proof line 5252) and remark divisors-remark-base-change-relative-assassin (proof line 5253); Same chapter flat.tex, definition labelled definition-pure and the discussion following it (proof lines 5262-5264). Turns purity into the statement that every relative associated point specializes into the fibre, describes the relative assassin after base changing to the henselization as an inverse image, and supplies the specialization criterion for purity that the converse must establish.
\item Associated points of the pullback along an etale morphism. Same chapter flat.tex, label lemma-etale-weak-assassin-up-down, invoked at proof line 5268 in the form applied to the morphism of fibres. Shows that a preimage of a relative associated point along the etale morphism g is again an associated point of the pulled back module on the corresponding fibre, which lets purity of the pullback be used to produce a specialization inside the fibre.
\end{itemize}\medskip
\ifdefined\hypertarget\hypertarget{sourcepointer-1}{}\fi
\noindent\textit{Source reference.} Original lemma 05J4: sources/stacks-flat/pinned-original-source.tex, source lines 4808--4808. Navigation pointer for \textbackslash{}ref\{definition-pure\} used inside the counted span (line 5264); the argument that consumes it is retained in the span.\par\medskip
\ifdefined\hypertarget\hypertarget{sourcepointer-2}{}\fi
\noindent\textit{Source reference.} Original lemma 05DC: sources/stacks-chow/tags.txt, source lines 6950--6950. Navigation pointer for \textbackslash{}ref\{divisors-lemma-base-change-relative-assassin\} used inside the counted span (line 5252); the argument that consumes it is retained in the span.\par\medskip
\ifdefined\hypertarget\hypertarget{sourcepointer-3}{}\fi
\noindent\textit{Source reference.} Original lemma 05KL: sources/stacks-chow/tags.txt, source lines 7205--7205. Navigation pointer for \textbackslash{}ref\{divisors-remark-base-change-relative-assassin\} used inside the counted span (line 5253); the argument that consumes it is retained in the span.\par\medskip
\ifdefined\hypertarget\hypertarget{sourcepointer-4}{}\fi
\noindent\textit{Source reference.} Original lemma 05L3: sources/stacks-flat/pinned-original-source.tex, source lines 5175--5175. Navigation pointer for \textbackslash{}ref\{lemma-associated-point-specializes\} used inside the counted span (line 5239); the argument that consumes it is retained in the span.\par\medskip
\ifdefined\hypertarget\hypertarget{sourcepointer-5}{}\fi
\noindent\textit{Source reference.} Original lemma 05FP: sources/stacks-flat/pinned-original-source.tex, source lines 353--353. Navigation pointer for \textbackslash{}ref\{lemma-etale-weak-assassin-up-down\} used inside the counted span (line 5268); the argument that consumes it is retained in the span.\par\medskip
\ifdefined\hypertarget\hypertarget{sourcepointer-6}{}\fi
\noindent\textit{Source reference.} Original lemma 05J9: sources/stacks-flat/pinned-original-source.tex, source lines 4949--4949. Navigation pointer for \textbackslash{}ref\{lemma-flat-descend-pure\} used inside the counted span (line 5228); the argument that consumes it is retained in the span.\par\medskip
\ifdefined\hypertarget\hypertarget{sourcepointer-7}{}\fi
\noindent\textit{Source reference.} Original lemma 05J7: sources/stacks-flat/pinned-original-source.tex, source lines 4908--4908. Navigation pointer for \textbackslash{}ref\{lemma-quasi-finite-base-change\} used inside the counted span (line 5226); the argument that consumes it is retained in the span.\par\medskip
\begin{lemma}
\label{lemma-criterion}
Let $f : X \to S$ be a morphism of schemes of finite type.
Let $\mathcal{F}$ be a quasi-coherent $\mathcal{O}_X$-module
of finite type. Let $s \in S$. Let $(S', s') \to (S, s)$ be an
elementary \'etale neighbourhood and let
$$
\xymatrix{
X \ar[d] & X' \ar[l]^g \ar[d] \\
S & S' \ar[l]
}
$$
be a commutative diagram of morphisms of schemes. Assume
\begin{enumerate}
\item $\mathcal{F}$ is flat over $S$ at all points of $X_s$,
\item $X' \to S'$ is of finite type,
\item $g^*\mathcal{F}$ is pure along $X'_{s'}$,
\item $g : X' \to X$ is \'etale, and
\item $g(X')$ contains $\text{Ass}_{X_s}(\mathcal{F}_s)$.
\end{enumerate}
In this situation $\mathcal{F}$ is pure along $X_s$ if and only
if the image of $X' \to X \times_S S'$ contains the points of
$\text{Ass}_{X \times_S S'/S'}(\mathcal{F} \times_S S')$
lying over points in $S'$ which specialize to $s'$.
\end{lemma}

\begin{proof}
Since the morphism $S' \to S$ is \'etale, we see that if $\mathcal{F}$
is pure along $X_s$, then $\mathcal{F} \times_S S'$ is pure along
$X_s$, see
Lemma \ref{lemma-quasi-finite-base-change}.
Since purity satisfies flat descent, see
Lemma \ref{lemma-flat-descend-pure},
we see that if $\mathcal{F} \times_S S'$ is pure along $X_{s'}$, then
$\mathcal{F}$ is pure along $X_s$. Hence we may replace $S$ by $S'$
and assume that $S = S'$ so that $g : X' \to X$ is an \'etale morphism
between schemes of finite type over $S$. Moreover, we may replace
$S$ by $\Spec(\mathcal{O}_{S, s})$ and assume that $S$ is local.

\medskip\noindent
First, assume that $\mathcal{F}$ is pure along $X_s$.
In this case every point of $\text{Ass}_{X/S}(\mathcal{F})$
specializes to a point of $X_s$ by purity. Hence by
Lemma \ref{lemma-associated-point-specializes}
we see that every point of $\text{Ass}_{X/S}(\mathcal{F})$
specializes to a point of $\text{Ass}_{X_s}(\mathcal{F}_s)$.
Thus every point of $\text{Ass}_{X/S}(\mathcal{F})$ is in the
image of $g$ (as the image is open and contains
$\text{Ass}_{X_s}(\mathcal{F}_s)$).

\medskip\noindent
Conversely, assume that $g(X')$ contains $\text{Ass}_{X/S}(\mathcal{F})$.
Let $S^h = \Spec(\mathcal{O}_{S, s}^h)$ be the henselization
of $S$ at $s$. Denote $g^h : (X')^h \to X^h$ the base change of $g$
by $S^h \to S$, and denote $\mathcal{F}^h$ the pullback of $\mathcal{F}$
to $X^h$. By
Divisors, Lemma \ref{divisors-lemma-base-change-relative-assassin} and
Remark \ref{divisors-remark-base-change-relative-assassin}
the relative assassin $\text{Ass}_{X^h/S^h}(\mathcal{F}^h)$
is the inverse image of $\text{Ass}_{X/S}(\mathcal{F})$ via the projection
$X^h \to X$. As we have assumed that $g(X')$ contains
$\text{Ass}_{X/S}(\mathcal{F})$ we conclude that the base change
$g^h((X')^h) = g(X') \times_S S^h$ contains
$\text{Ass}_{X^h/S^h}(\mathcal{F}^h)$. In this way
we reduce to the case where $S$ is the spectrum of a henselian local ring.
Let $x \in \text{Ass}_{X/S}(\mathcal{F})$. To finish the proof of the
lemma we have to show that $x$ specializes to a point of $X_s$, see
criterion (4) for purity in discussion following
Definition \ref{definition-pure}.
By assumption there exists a $x' \in X'$ such that $g(x') = x$.
As $g : X' \to X$ is \'etale, we see that
$x' \in \text{Ass}_{X'/S}(g^*\mathcal{F})$, see
Lemma \ref{lemma-etale-weak-assassin-up-down} (applied to
the morphism of fibres $X'_w \to X_w$ where $w \in S$ is the image of $x'$).
Since $g^*\mathcal{F}$ is pure along $X'_s$ we see that $x' \leadsto y$
for some $y \in X'_s$. Hence $x = g(x') \leadsto g(y)$ and
$g(y) \in X_s$ as desired.
\end{proof}

\end{document}
